\documentclass{situationist-preprint}

\title{Counting Streams of Consciousness}
\msid{MND-15-b}
\msfamily{MND-15}
\msdiscipline{Philosophy of Mind}
\mstype{Companion manuscript}
\mskeywords{unity of consciousness; streams; split brain; partial unity; co-consciousness}

\begin{document}
\maketitle

\begin{abstract}
How many streams of consciousness does a subject have at a time? The question is usually answered by counting total experiences, which yields two in the split-brain patient and one in the normal subject. We set out the method of counting that the account of partial unity in MND-15-a induces: streams are the connected components of the co-consciousness relation, and total experiences are its maximal cliques. The two counts coincide when co-consciousness is transitive and diverge when it is not. On the proposed method the split-brain patient has one stream with two maximal experiences, which is what partial unity says.
\end{abstract}

\section{Two ways of counting}

MND-15-a defines co-consciousness as a relation between simultaneous experiences of a subject that are parts of a single experience (Definition~2.1), and a stream at a time as a maximal set of experiences linked by chains of co-conscious pairs (Definition~2.2). Represent a subject's experiences at a time as the vertices of a graph, joined when they are co-conscious. Two counts are then available.

\begin{definition}[Two counts]\label{d:counts}
The \emph{stream count} is the number of connected components of the graph. The \emph{total-experience count} is the number of maximal sets of pairwise co-conscious experiences, its maximal cliques; each such set composes one total experience, something it is like to have all of its members together \citep{BayneChalmers2003}.
\end{definition}

\begin{proposition}\label{p:coincide}
The two counts coincide if and only if every component of the graph is a clique, that is, if and only if co-consciousness is transitive among the subject's experiences at that time.
\end{proposition}

\begin{proof}
Every clique lies within one component, and every component contains at least one maximal clique. If each component is a clique, it is its own unique maximal clique, so the counts are equal. If some component is not a clique, it contains two experiences that are not co-conscious, each of which lies in some maximal clique; no clique contains both, so that component contains at least two maximal cliques and the total-experience count exceeds the stream count.
\end{proof}

Traditional discussions count total experiences and call the result the number of streams \citep{Nagel1971,Bayne2010}. Proposition~\ref{p:coincide} shows that this is harmless under transitivity, which is the normal case established in Lemma~3.1 of MND-15-a, and misleading otherwise.

\section{The split-brain case}

In the split-brain patient, by the argument of MND-15-a (Lemma~4.1 and Theorem~4.3), experiences of the left visual field and of the right visual field are not co-conscious, but each is co-conscious with shared bodily and emotional experiences, which are integrated with both by Principle~P1 \citep{Pinto2017}. The graph is connected and not a clique. The stream count is one; the total-experience count is at least two. The patient has one stream, as MND-15-a holds, containing two overlapping total experiences.

\begin{objection}[Only total experiences matter]
What it is like for the subject is fixed by her total experiences, so their number is the number that matters.
\end{objection}

\begin{reply}
It fixes what it is like within each total experience, but it does not settle how many streams there are, since a stream is individuated by what can be reached by co-consciousness, and the overlapping parts of the two total experiences are not had twice. Principle~P2 of MND-15-a forbids duplicating the shared bodily experience into two tokens, and without duplication the two total experiences share parts and cannot belong to separate streams \citep{Hurley1998,Lockwood1989}.
\end{reply}

The method applies wherever integration is incomplete. It is the method MND-23 uses for the octopus, whose arm-related contents may form a connected but non-transitive graph: one stream, possibly several maximal experiences.

\section{Conclusion}

Streams are counted by components and total experiences by maximal cliques. The counts agree in normal subjects and diverge under partial unity, and the split-brain patient has one partially unified stream.

\begin{thebibliography}{9}\small
\bibitem[Bayne(2010)]{Bayne2010} Bayne, T. (2010). \emph{The Unity of Consciousness}. Oxford University Press.
\bibitem[Bayne and Chalmers(2003)]{BayneChalmers2003} Bayne, T. and Chalmers, D.~J. (2003). What is the unity of consciousness? In A. Cleeremans (ed.), \emph{The Unity of Consciousness: Binding, Integration, and Dissociation}. Oxford University Press.
\bibitem[Hurley(1998)]{Hurley1998} Hurley, S.~L. (1998). \emph{Consciousness in Action}. Harvard University Press.
\bibitem[Lockwood(1989)]{Lockwood1989} Lockwood, M. (1989). \emph{Mind, Brain and the Quantum: The Compound \enquote{I}}. Blackwell.
\bibitem[Nagel(1971)]{Nagel1971} Nagel, T. (1971). Brain bisection and the unity of consciousness. \emph{Synthese}, 22(3--4), 396--413.
\bibitem[Pinto et~al.(2017)]{Pinto2017} Pinto, Y., Neville, D.~A., Otten, M., Corballis, P.~M., Lamme, V.~A.~F., de~Haan, E.~H.~F., Foschi, N. and Fabri, M. (2017). Split brain: Divided perception but undivided consciousness. \emph{Brain}, 140(5), 1231--1237.
\end{thebibliography}

\end{document}
